% countdown-conops.tex
% Concept of Operations for the countdown program.
% Structure per IEEE Std 1362-1998 (ConOps) / ISO/IEC/IEEE 29148:2018.
% Compiled with: pdflatex countdown-conops.tex

\documentclass[12pt]{article}

\usepackage{geometry}
\geometry{letterpaper, margin=1in}
\usepackage{hyperref}
\usepackage{booktabs}
\usepackage{array}
\usepackage{longtable}
\usepackage{enumitem}
\usepackage{titlesec}
\usepackage{fancyhdr}
\usepackage{xcolor}

\definecolor{cdteal}{RGB}{0,90,120}

\pagestyle{fancy}
\fancyhf{}
\renewcommand{\headrulewidth}{0.4pt}
\fancyhead[L]{\small Countdown Project}
\fancyhead[R]{\small ConOps: countdown}
\fancyfoot[C]{\thepage}

\titleformat{\section}{\large\bfseries\color{cdteal}}{\thesection}{0.5em}{}[\titlerule]
\titleformat{\subsection}{\normalsize\bfseries}{\thesubsection}{0.5em}{}

\newcommand{\code}[1]{\texttt{#1}}
\newcommand{\req}[1]{\texttt{#1}}

\begin{document}

\begin{titlepage}
  \centering
  \vspace*{2cm}
  {\Huge\bfseries\color{cdteal} Concept of Operations\\[0.5em]}
  {\Large countdown: An X~Window Countdown Timer}\\[2em]
  {\large Countdown Software Project}\\[0.5em]
  {\normalsize Prepared for: Users and Maintainers of \code{countdown}}\\[0.5em]
  {\normalsize Prepared by: J.\,R.\,Evans}\\[2em]
  \begin{tabular}{ll}
    Document identifier: & CD-CONOPS-001                                  \\
    Revision:            & 1.0                                            \\
    Date:                & June 13, 2026                                  \\
    Status:              & Released                                       \\
    Governing standard:  & IEEE Std 1362-1998 (ConOps)                    \\
  \end{tabular}
  \vfill
  {\small\textit{This document is prepared in accordance with the
  principles of rational documentation articulated in
  D.\,L.\,Parnas and P.\,C.\,Clements, ``A Rational Design Process:
  How and Why to Fake It,'' \textit{IEEE Transactions on Software
  Engineering}, vol.\,SE-12, no.\,2, pp.\,251--257, February 1986.}}
\end{titlepage}

\tableofcontents
\clearpage

% =================================================================
\section{Scope}
% =================================================================

\subsection{System identification}

This Concept of Operations (ConOps) describes \code{countdown}, a small
single-window graphical application for the X~Window System.  Given a
fixed future instant --- by default \code{2026-08-19 08:00:00} in local
time --- the program displays the time remaining until that instant and
updates the display live until it is reached.

\subsection{Document purpose}

The ConOps states, in the language of the user, what \code{countdown} is
for, how the situation looks without it, why it is worth building, and how
a user interacts with it.  It is the topmost document in the
\code{countdown} documentation set; the Software Requirements
Specification (CD-SRS-001), Software Design Document (CD-SDD-001), and
Software Test Plan (CD-STP-001) refine and verify the intent recorded
here.

\subsection{A note on rational documentation}

Following Parnas and Clements, this document is written as though the
concept had been arrived at by an orderly process.  The real path was the
usual mixture of trial and revision; the value of presenting an idealised
account is that a reader can understand the system without retracing every
false start.  The same fiction is maintained, deliberately, throughout the
documentation set.

% =================================================================
\section{Referenced documents}
% =================================================================

\begin{longtable}{p{0.28\textwidth}p{0.66\textwidth}}
\toprule
\textbf{Reference} & \textbf{Title} \\
\midrule
IEEE Std 1362-1998 & Guide for Information Technology---System Definition---Concept of Operations Document \\
ISO/IEC/IEEE 29148:2018 & Systems and software engineering---Life cycle processes---Requirements engineering \\
Parnas \& Clements, 1986 & A Rational Design Process: How and Why to Fake It \\
Knuth, 1984 & Literate Programming, \textit{The Computer Journal} 27(2) \\
CD-SRS-001 & Software Requirements Specification: countdown \\
CD-SDD-001 & Software Design Document: countdown \\
CD-STP-001 & Software Test Plan: countdown \\
X11R7 / Xlib & \textit{Xlib---C Language X Interface}, X Consortium \\
POSIX.1-2017 & IEEE Std 1003.1-2017 (\code{time}, \code{mktime}, \code{select}) \\
\bottomrule
\end{longtable}

% =================================================================
\section{Current system or situation}
% =================================================================

\subsection{Background}

People frequently wish to mark the approach of a specific moment: the
start of an event, a deadline, a launch, a celebration.  In the absence of
a dedicated tool the countdown is performed by mental arithmetic against a
wall clock or calendar, or by a web page of uncertain provenance that
requires a network connection and a browser.

\subsection{Description of the current situation}

On a typical desktop running the X~Window System the user today has:

\begin{itemize}[noitemsep]
  \item the \code{date} command, which prints the current time but
        performs no subtraction toward a target;
  \item general-purpose calculators, into which the user must transcribe
        and difference two timestamps by hand;
  \item online countdown widgets, which depend on a browser, a network
        path, and a third party's continued operation.
\end{itemize}

\subsection{Problems and limitations}

The manual approach is error-prone (time-zone and daylight-saving
mistakes are common), is not continuously updated, and gives no glanceable
visual artefact that can be left running on the desktop.  The online
approach surrenders privacy and availability to a remote service for a
computation the local machine can trivially perform.

% =================================================================
\section{Justification for change}
% =================================================================

A countdown to a known instant is a self-contained computation requiring
only the system clock and the local time-zone rules.  A small native
program can perform it correctly, continuously, and offline, presenting a
single uncluttered window.  The justification for building
\code{countdown} is therefore:

\begin{itemize}[noitemsep]
  \item \textbf{Correctness:} time-zone and daylight-saving handling is
        delegated to the operating system's \code{mktime}, eliminating a
        common class of manual error.
  \item \textbf{Continuity:} the display advances automatically, without
        user action, at sub-second resolution.
  \item \textbf{Independence:} no network, browser, or third-party service
        is involved; the program runs wholly on the local host.
  \item \textbf{Transparency:} the program is a literate document, so its
        behaviour can be read and audited in full.
\end{itemize}

% =================================================================
\section{Proposed system}
% =================================================================

\subsection{Overview}

\code{countdown} is a single executable that opens one window on the
X~Window System.  The window shows the target instant, the remaining time
expressed as days and a \code{HH:MM:SS} clock, a status line, and a one-line
hint listing the available keys.  The user starts and pauses the countdown
from the keyboard or mouse and quits when finished.

\subsection{Operational environment}

The program runs on any POSIX host providing an X~server (Linux with an
X~display, or macOS with XQuartz).  It links only against the core Xlib
client library; no widget toolkit is required.  It consumes negligible CPU
because it sleeps on \code{select} between updates rather than polling.

\subsection{High-level capabilities}

\begin{enumerate}[noitemsep]
  \item Count down to a default instant of \code{2026-08-19 08:00:00},
        or to any instant supplied on the command line.
  \item Display the remaining time live, updating at least once per
        second.
  \item Let the user start, pause, resume, and quit interactively.
  \item Announce arrival when the target instant is reached, and never
        display a negative remaining time.
\end{enumerate}

\subsection{Modes of operation}

The program has three user-visible modes: \emph{armed} (started but not
yet running --- the initial state, showing the full duration);
\emph{running} (counting down live); and \emph{paused} (the figure is
frozen at the value captured when the user paused).  A fourth, terminal
condition, \emph{arrived}, is entered when the remaining time reaches zero.

% =================================================================
\section{Operational scenarios}
% =================================================================

\subsection{Scenario 1: Counting down to the default instant}

Dana wants a visible countdown to the morning of 19~August 2026 left
running on her desktop.  She launches \code{countdown} with no arguments.
A window appears showing \emph{Target: Wed 19 Aug 2026 08:00:00} and the
full remaining duration with the status \emph{ready}.  She presses the
space bar; the status changes to \emph{running} and the seconds begin to
tick down.  She leaves the window in a corner of her screen.

\subsection{Scenario 2: A custom target}

Raj is timing a 25-minute presentation rehearsal that should end at
14:30.  He runs \code{countdown "2026-06-13 14:30:00"}.  The window opens
with that instant as the target.  Part-way through he is interrupted; he
clicks the window to pause, deals with the interruption, then clicks again
to resume.

\subsection{Scenario 3: Reaching the target}

As the target instant arrives, the remaining figure reaches
\code{0 d 00:00:00} and the status line changes to \emph{THE MOMENT HAS
ARRIVED}.  The figure does not go negative.  The user observes the message
and presses \code{q} to quit.

\subsection{Scenario 4: A mistyped target}

A user runs \code{countdown "Aug 19"}.  Because the argument does not match
\code{YYYY-MM-DD HH:MM:SS}, the program prints a diagnostic naming the
expected form and exits with a non-zero status, without opening a window.

% =================================================================
\section{Summary of impacts}
% =================================================================

\subsection{Operational impacts}

Users gain a reliable, glanceable countdown that runs locally.  No new
infrastructure, account, or network dependency is introduced.

\subsection{Organisational impacts}

None of consequence.  The program is a single self-contained executable;
there is no server to administer and no data to retain between runs.

\subsection{Support and maintenance impacts}

Because the program is a literate document, maintenance proceeds by
reading and editing \code{countdown.w} and re-tangling.  A maintainer
needs familiarity with C, Xlib, and the CWEB tools, all of which are
documented in the program itself and in CD-SDD-001.

% =================================================================
\section{Analysis of the proposed system}
% =================================================================

\subsection{Alternatives considered}

\textbf{A terminal (text) countdown.}  Simplest to build, but the ConOps
calls explicitly for a window on the X~Window System, and a graphical
artefact is more glanceable.  Rejected on requirement grounds.

\medskip\noindent
\textbf{A toolkit-based GUI (GTK, Qt, Motif).}  Would give buttons and
theming for free, but pulls in a large dependency for a one-window program
and obscures the underlying mechanism.  Rejected in favour of raw Xlib,
which keeps the program small, dependency-light, and fully readable.

\medskip\noindent
\textbf{A browser/web widget.}  Rejected for the privacy, availability,
and dependency reasons set out in Section~4.

\subsection{Trade-offs}

Choosing raw Xlib trades developer convenience (no ready-made widgets) for
a minimal, transparent, portable program.  Choosing \code{select}-based
multiplexing over a busy-wait trades a few lines of code for near-zero idle
CPU.  Choosing local-time semantics (via \code{mktime}) trades a little
ambiguity at daylight-saving boundaries for the convenience of letting
users write wall-clock times.

\subsection{Justification for the selected approach}

The selected approach---a single Xlib window driven by a
\code{select}-based event-and-timer loop, written as a literate
program---meets every capability in Section~6 with the least machinery and
the greatest readability, and is therefore the most defensible choice for
a tool whose whole purpose is to be small, correct, and self-evident.

\end{document}
